\section{Operational Definitions}

\begin{description}[leftmargin=3cm, style=nextline]

\item[Reliability]
Three measurable properties: (1) \textit{job success rate} --- the percentage of
\dagster{} runs completing without error at each concurrency level; (2) \textit{failure
isolation} --- whether one failing job measurably affects other concurrently running jobs;
(3) \textit{\mttr{}} --- time in seconds from failure detection to successful job
re-execution.

\item[System Stability]
Measured primarily through \textit{execution time variance} --- the standard deviation of
job execution time across repeated runs at each concurrency level. Variance reveals whether
system behaviour is consistent under load, which mean performance metrics cannot capture.
A system with high variance is operationally unstable even if its mean execution time is
acceptable.

\item[Scalability]
Measured as \textit{concurrency capacity} --- the maximum number of concurrent jobs the
system sustains while maintaining a job success rate above 95\%.

\item[Operational Overhead]
Three measurable components: (1) \textit{pod scheduling latency} --- time from \dagster{}
job submission to pod entering running state; (2) \textit{container startup time} --- time
from pod running state to first job instruction executing; (3) \textit{net execution time
delta} --- difference in total end-to-end execution time between \vm{} and Kubernetes for
identical jobs at identical concurrency levels.

\item[Crossover Point]
Formally defined as the concurrent workload level at which \textbf{(1)} the \vm{} job
success rate drops below 95\% due to resource contention, \textbf{AND (2)} the Kubernetes
total execution time including overhead becomes lower than the \vm{} execution time under
active contention. This is a computable result derived directly from experimental data, not
an interpretation.

\end{description}
